\section{What information defines a line}

A line is one of the simplest geometric objects, but it is also one of the easiest objects to misunderstand. In elementary algebra a line is often introduced through an equation such as \(y=mx+b\). That equation is useful, but it should not be confused with the object itself. A line is not a formula. A line is a set of points arranged in one direction.

Analytic geometry studies this set of points by translating geometric information into algebraic language. The first question is therefore not ``Which formula should I use?'' The first question is: what information is enough to determine the line?

\subsection*{A point and a direction}

The most natural description of a line uses two pieces of information:
\begin{itemize}
\item one point on the line,
\item one non-zero vector that gives the direction of the line.
\end{itemize}

If \(A\) is a point and \(v\ne0\) is a direction vector, then the line consists of all points obtained by starting at \(A\) and moving by arbitrary multiples of \(v\).

\begin{definitionbox}
Let \(A\) be a point and let \(v\ne0\) be a vector. The line through \(A\) in direction \(v\) is the set of points
\[
P=A+t v,
\]
where \(t\) runs through all real numbers.
\end{definitionbox}

The parameter \(t\) controls the movement along the line. When \(t=0\), we are at the point \(A\). When \(t=1\), we move from \(A\) by exactly one copy of the vector \(v\). When \(t=-1\), we move by the same vector in the opposite direction.

\begin{examplebox}
Let
\[
A=(1,2),
\qquad
v=(3,-1).
\]
The line through \(A\) in direction \(v\) is
\[
P=(1,2)+t(3,-1).
\]
For \(t=0\), we get
\[
P=(1,2).
\]
For \(t=1\), we get
\[
P=(4,1).
\]
For \(t=-1\), we get
\[
P=(-2,3).
\]
These are different points, but they all lie on the same line because they are obtained from the same starting point by moving in the same direction.
\end{examplebox}

This description matches the language developed in the previous chapter. A point gives position. A vector gives direction. A parameter tells us how much of that direction to use.

\subsection*{Two distinct points}

A second common way to determine a line is to give two distinct points. If the points are different, then the vector from one point to the other gives a direction.

Let
\[
A=(a_1,a_2),
\qquad
B=(b_1,b_2),
\qquad
A\ne B.
\]
Then
\[
\overrightarrow{AB}=B-A=(b_1-a_1,\,b_2-a_2)
\]
is a non-zero direction vector for the line through \(A\) and \(B\). Therefore the line can be written as
\[
P=A+t(B-A).
\]

The condition \(A\ne B\) is not a technical decoration. If \(A=B\), then the vector \(B-A\) is the zero vector, and the zero vector has no direction. One point alone does not determine a unique line.

\begin{examplebox}
Let
\[
A=(2,-1),
\qquad
B=(5,3).
\]
Then
\[
B-A=(3,4).
\]
The line through \(A\) and \(B\) can be written as
\[
P=(2,-1)+t(3,4).
\]
For \(t=1\), this gives \((5,3)\), so the point \(B\) is indeed on the line.
\end{examplebox}

\subsection*{A point and a normal vector in the plane}

In the plane there is another important way to determine a line. Instead of giving a direction along the line, we may give a direction perpendicular to the line. Such a vector is called a normal vector.

If \(n\ne0\) is perpendicular to a line and \(A\) is a point on the line, then a point \(P\) lies on the line exactly when the displacement \(P-A\) is perpendicular to \(n\). Using the dot product, this becomes
\[
n\cdot(P-A)=0.
\]

This form is the beginning of the normal description of a line. It is extremely useful because perpendicularity, distance from a point to a line, and the general equation of a line all become natural from it.

\begin{remarkbox}
A direction vector points along the line. A normal vector points across the line. Both can describe the same line, but they reveal different information.
\end{remarkbox}

\subsection*{The same object, different data}

A line may therefore be determined from different kinds of information. We may know a point and a direction. We may know two points. In the plane, we may know a point and a normal vector. We may also know an equation and need to read from it what the point, direction, or normal vector is.

This is why this chapter will not present only one formula for a line. A line has many algebraic descriptions. Each description is a way of reading the same geometric object.

\begin{summarybox}
When first meeting a line problem, ask:
\begin{itemize}
\item Am I given a point on the line?
\item Am I given a direction vector?
\item Am I given two points?
\item Am I given a normal vector or an implicit equation?
\item Is the line in the plane or in space?
\item Which description makes the given information visible?
\end{itemize}
\end{summarybox}

The goal is not to memorize all possible formulas at once. The goal is to understand what information each formula is carrying.